% GENERATED FILE. Edit canonical lessons, then run npm run generate:books.
% canonical-source-hash: fnv1a64:a472d42ffbab22f2
% canonical-lessons: JA-C253-futte-imasu, JA-C253-toki, JA-C253-nenjuu, JA-W253-nen-kanji, JA-C253-kawa-no-hi, JA-R253-what-was-going-on

\chapter{What Was Going On}
\label{ch:ja-what-was-going-on}
\hlchaptermodality{253}

\begin{chapteropening}
\textbf{By the end of this chapter:} \emph{I can say what is and was going on (ame ga futte imashita), set a scene with toki (kaetta toki, haha wa nete imashita), and follow a short story of a day at the river.}

Complete JA-R253-what-was-going-on: set the scene of a story and say what happened in it.
\end{chapteropening}

\section[\texorpdfstring{\ja{ふっています} … (futte imasu …)}{ふっています … (futte imasu …)}]{\ja{ふっています} / \ja{ふっていました} (futte imasu / futte imashita) — what is going on}

\label{lesson:JA-C253-futte-imasu}



\begin{quote}
Say \textquotedblleft{}It will probably rain tomorrow\textquotedblright{}. (\emph{ashita ame ga furu deshō}.) Then the \textbf{-\ja{て}} form of \textbf{\ja{ふる}}. (\emph{futte}.)
\end{quote}

\subsection*{You'll want to know: -\ja{ています}}
\begin{quote}
\textbf{\ja{あめが} \ja{ふっています。}}
\end{quote}

\begin{quote}
\emph{ame ga futte imasu.} — \textquotedblleft{}It is raining.\textquotedblright{}
\end{quote}

The \textbf{-\ja{て}} form with \textbf{\ja{います}} (\emph{imasu}), \textquotedblleft{}is there\textquotedblright{}, says what is going on: the rain has started and has not stopped.

\begin{quote}
\textbf{\ja{いま} \ja{あめが} \ja{ふっています。}}
\end{quote}

\begin{quote}
\emph{ima ame ga futte imasu.} — \textquotedblleft{}It is raining now.\textquotedblright{}
\end{quote}

\begin{quote}
\textbf{\ja{ともだちを} \ja{まっています。}}
\end{quote}

\begin{quote}
\emph{tomodachi o matte imasu.} — \textquotedblleft{}I am waiting for a friend.\textquotedblright{}
\end{quote}

\subsection*{You'll want to know: -\ja{ていました}}
With \textbf{\ja{いました}} (\emph{imashita}), it was going on:

\begin{quote}
\textbf{\ja{きのう} \ja{あめが} \ja{ふっていました。}}
\end{quote}

\begin{quote}
\emph{kinō ame ga futte imashita.} — \textquotedblleft{}It was raining yesterday.\textquotedblright{}
\end{quote}

\begin{quote}
\textbf{\ja{ははは} \ja{しんぶんを} \ja{よんでいました。}}
\end{quote}

\begin{quote}
\emph{haha wa shinbun o yonde imashita.} — \textquotedblleft{}My mother was reading the newspaper.\textquotedblright{}
\end{quote}

A story needs both kinds of past. \textbf{-\ja{ていました}} is the scene, going on; \textbf{-\ja{ました}} is the event, once, inside it:

\begin{quote}
\textbf{\ja{あめが} \ja{ふっていました。}}
\end{quote}

\begin{quote}
\emph{ame ga futte imashita.} — \textquotedblleft{}It was raining.\textquotedblright{}
\end{quote}

\begin{quote}
\textbf{\ja{でんしゃが} \ja{とまりました。}}
\end{quote}

\begin{quote}
\emph{densha ga tomarimashita.} — \textquotedblleft{}The train stopped.\textquotedblright{}
\end{quote}

\subsection*{Your turn}
Say these aloud:

\begin{itemize}
\raggedright
  \item \textquotedblleft{}My mother was reading the newspaper\textquotedblright{} — \emph{haha wa shinbun o yonde imashita}
  \item \textquotedblleft{}My father was reading the newspaper\textquotedblright{} — \emph{chichi wa shinbun o yonde imashita}
  \item \textquotedblleft{}My elder sister was reading the newspaper\textquotedblright{} — \emph{ane wa shinbun o yonde imashita}
\end{itemize}

\begin{tcolorbox}[breakable,colback=teal!4,colframe=teal!35!black,title={Before you move on}]
\textquotedblleft{}It is raining\textquotedblright{}? (\textbf{\ja{あめが} \ja{ふっています}}.) \textquotedblleft{}It was raining\textquotedblright{}? (\textbf{\ja{ふっていました}}.) \textquotedblleft{}I am waiting\textquotedblright{}? (\textbf{\ja{まっています}}.) Which is the scene, \textbf{\ja{ふっていました}} or \textbf{\ja{とまりました}}? (\textbf{\ja{ふっていました}}.)
\end{tcolorbox}
\section[\texorpdfstring{\ja{とき} (toki)}{とき (toki)}]{\ja{とき} (toki) — when; the time when}

\label{lesson:JA-C253-toki}



\begin{quote}
Before the new word: say \textquotedblleft{}My mother was asleep\textquotedblright{}. (\emph{haha wa nete imashita}.) Then say \textquotedblleft{}I got up\textquotedblright{}, to a friend. (\emph{okita}.)
\end{quote}

\subsection*{You'll want to know: \ja{とき}}
\textbf{\ja{とき}} — \emph{toki} — \textquotedblleft{}when\textquotedblright{}, \textquotedblleft{}the time when\textquotedblright{}. The plain form goes before it, and \textbf{\ja{、}} closes the time:

\begin{quote}
\textbf{\ja{かえった} \ja{とき、ははは} \ja{ねていました。}}
\end{quote}

\begin{quote}
\emph{kaetta toki, haha wa nete imashita.} — \textquotedblleft{}When I got home, my mother was asleep.\textquotedblright{}
\end{quote}

\begin{quote}
\textbf{\ja{おきた} \ja{とき、あめが} \ja{ふっていました。}}
\end{quote}

\begin{quote}
\emph{okita toki, ame ga futte imashita.} — \textquotedblleft{}When I got up, it was raining.\textquotedblright{}
\end{quote}

After a noun it takes \textbf{\ja{の}}:

\begin{quote}
\textbf{\ja{こどもの} \ja{とき、うみに} \ja{いきました。}}
\end{quote}

\begin{quote}
\emph{kodomo no toki, umi ni ikimashita.} — \textquotedblleft{}When I was a child, I went to the sea.\textquotedblright{}
\end{quote}

\begin{cousinweb}
\textbf{\ja{とき}} (\emph{toki}) is a native Japanese word (wago), inherited rather than borrowed from Chinese. JMdict writes it with one kanji and glosses it \textquotedblleft{}time; hour; moment\textquotedblright{} and \textquotedblleft{}occasion; case\textquotedblright{}; KANJIDIC2 glosses that kanji \textquotedblleft{}time, hour\textquotedblright{} and gives \ja{とき} as its native reading. It is the \textbf{\ja{とき}} inside \textbf{\ja{ときどき}}, \textquotedblleft{}sometimes\textquotedblright{} (JMdict, KANJIDIC2).
\end{cousinweb}

\subsection*{Your turn}
Say these aloud:

\begin{itemize}
\raggedright
  \item \emph{toki}
  \item \textquotedblleft{}When I got home, my mother was asleep\textquotedblright{} — \emph{kaetta toki, haha wa nete imashita}
  \item \textquotedblleft{}When I got home, my father was asleep\textquotedblright{} — \emph{kaetta toki, chichi wa nete imashita}
\end{itemize}

\begin{tcolorbox}[breakable,colback=teal!4,colframe=teal!35!black,title={Before you move on}]
\textquotedblleft{}When I was a child\textquotedblright{}? (\textbf{\ja{こどもの} \ja{とき}}.) \textquotedblleft{}When I got home\textquotedblright{}? (\textbf{\ja{かえった} \ja{とき}}.)
\end{tcolorbox}
\section[\texorpdfstring{\ja{年}\ja{じゅう} (nenjū)}{年じゅう (nenjū)}]{\ja{年}\ja{じゅう} (nenjū) — all year round}

\label{lesson:JA-C253-nenjuu}



\begin{quote}
Before the new word: say \textquotedblleft{}when I was a child\textquotedblright{}. (\emph{kodomo no toki}.) Then say \textquotedblleft{}hot\textquotedblright{}. (\emph{atsui}.)
\end{quote}

\subsection*{You'll want to know: \ja{年}\ja{じゅう}}
\textbf{\ja{年}\ja{じゅう}} — \emph{nenjū} — \textquotedblleft{}all year round\textquotedblright{}, the whole year through.

\begin{quote}
\textbf{\ja{ここは} \ja{年}\ja{じゅう} \ja{あついです。}}
\end{quote}

\begin{quote}
\emph{koko wa nenjū atsui desu.} — \textquotedblleft{}It is hot here all year round.\textquotedblright{}
\end{quote}

Four beats: \emph{ne–n–ju–u}, the long \emph{ū} held for two. The first sign is new, read \textbf{\ja{ねん}} (\emph{nen}): the \emph{nen} of \textbf{\ja{らいねん}} and \textbf{\ja{きょねん}}. The next lesson writes it.

\begin{cousinweb}
\textbf{\ja{年}\ja{じゅう}} (\emph{nenjū}) is Sino-Japanese (kango): two Chinese characters read with their on'yomi, readings borrowed from Chinese. They are \emph{nen}, \textquotedblleft{}year\textquotedblright{} (Middle Chinese \emph{nen}) and \emph{jū}, \textquotedblleft{}in, inside, middle\textquotedblright{} (\emph{trjuwng}) (KANJIDIC2; Middle Chinese in Baxter's transcription). KANJIDIC2 gives that second kanji the on'yomi \emph{chū}; JMdict lists \emph{nenchū} as another reading of this same word, so \emph{jū} is \emph{chū} voiced. JMdict spells the word with two kanji, or with the second part in kana, as here, and glosses it \textquotedblleft{}all year round; throughout the year\textquotedblright{} (JMdict).
\end{cousinweb}

\subsection*{Your turn}
Say these aloud:

\begin{itemize}
\raggedright
  \item \emph{nenjū}
  \item \textquotedblleft{}It is hot here all year round\textquotedblright{} — \emph{koko wa nenjū atsui desu}
  \item \textquotedblleft{}It is cold here all year round\textquotedblright{} — \emph{koko wa nenjū samui desu}
\end{itemize}

\begin{tcolorbox}[breakable,colback=teal!4,colframe=teal!35!black,title={Before you move on}]
\textquotedblleft{}All year round\textquotedblright{}? (\textbf{\ja{年}\ja{じゅう}}, \emph{nenjū}.) Its first sign is read? (\textbf{\ja{ねん}}, \emph{nen}, as in \textbf{\ja{らいねん}}.)
\end{tcolorbox}
\section[\texorpdfstring{\ja{年} (nen)}{年 (nen)}]{\ja{年} — the year, six strokes}

\label{lesson:JA-W253-nen-kanji}



\begin{quote}
Say \emph{nenjū}, all year round. Then call up the four strokes of \textbf{\ja{月}}: left side, top and right side, two bars.
\end{quote}

\subsection*{Script — the shape}
\begin{quote}
\textbf{\ja{年}} — \emph{nen}
\end{quote}

Its parts:

\begin{itemize}
\raggedright
  \item a short falling stroke at the top
  \item three bars
  \item a short stem on the left
  \item a long centre stem
\end{itemize}

Six strokes, five lifts:

\begin{enumerate}
\raggedright
  \item draw the short falling stroke at the top, down to the left
  \item lift, draw the top bar from left to right
  \item lift, draw the second bar from left to right
  \item lift, draw the short stem on the left, down from the second bar
  \item lift, draw the long third bar from left to right
  \item lift, draw the long centre stem from the top bar straight down through the others
\end{enumerate}

The long stem comes last and runs through the bars below the top one. Here it is read \textbf{\ja{ねん}} (\emph{nen}), as in \textbf{\ja{年}\ja{じゅう}}; this lesson asks only for that reading.

\subsection*{Your turn}
\begin{enumerate}
\raggedright
  \item Find \textbf{\ja{年}} in \textbf{\ja{年}\ja{じゅう}} and say \emph{nenjū}.
  \item Copy \textbf{\ja{年}} three times, saying \emph{nen} at each one.
  \item Find it once more, in \textbf{\ja{ここは} \ja{年}\ja{じゅう} \ja{あついです}} (\emph{koko wa nenjū atsui desu}), and say the line.
\end{enumerate}

\begin{tcolorbox}[breakable,colback=teal!4,colframe=teal!35!black,title={Before you move on}]
How many strokes in all? (\textbf{Six}.)

Source: \href{https://github.com/KanjiVG/kanjivg/blob/master/kanji/05e74.svg}{KanjiVG stroke-order data}: KanjiVG, kanji/05e74.svg, kanji \ja{年} (U+5E74), 6 stroke paths kvg:05e74-s1 to s6 (Ulrich Apel and contributors, CC BY-SA 3.0). Only the order and direction are taken from it.
\end{tcolorbox}
\section[(reading)]{(reading) — a day at the river}

\label{lesson:JA-C253-kawa-no-hi}



\begin{quote}
Thirteen short lines, a summer day told in order. Every word in them is one you know.
\end{quote}

\begin{tcolorbox}[breakable,skin=enhanced,colback=orange!6,colframe=orange!45!black,boxrule=0.5pt,arc=1mm,left=6pt,right=6pt,top=4pt,bottom=4pt,fonttitle=\bfseries,title={Read this}]
\begin{quote}
\ja{なつやすみに} \ja{ちちと} \ja{かわに} \ja{いきました。} \ja{あさ} \ja{はやく} \ja{おきて、でかけました。} \ja{でんしゃに} \ja{のって、えきで} \ja{おりました。} \ja{それから、すこし} \ja{あるきました。} \ja{ついた} \ja{とき、あめが} \ja{ふっていました。} \ja{あめは} \ja{ひるに} \ja{やみました。} \ja{ちちと} \ja{かわで} \ja{およぎました。} \ja{みずは} \ja{きれいでした。} \ja{ひるごはんを} \ja{たべてから、さかなを} \ja{みました。} \ja{ちちは} \ja{しんぶんを} \ja{よんでいました。} \ja{そして、ゆうがた} \ja{うちに} \ja{かえりました。} \ja{かわの} \ja{みずは} \ja{年}\ja{じゅう} \ja{つめたいです。} \ja{月日}\ja{は} \ja{はやいです。}
\end{quote}

\emph{Read it:} the thirteen lines once, without stopping

Line by line:

\begin{enumerate}
\raggedright
  \item \textbf{\ja{なつやすみに} \ja{ちちと} \ja{かわに} \ja{いきました。}} (\emph{natsuyasumi ni chichi to kawa ni ikimashita.}) — \textquotedblleft{}In the summer holidays I went to the river with my father.\textquotedblright{}
  \item \textbf{\ja{あさ} \ja{はやく} \ja{おきて、でかけました。}} (\emph{asa hayaku okite, dekakemashita.}) — \textquotedblleft{}I got up early in the morning and set off.\textquotedblright{}
  \item \textbf{\ja{でんしゃに} \ja{のって、えきで} \ja{おりました。}} (\emph{densha ni notte, eki de orimashita.}) — \textquotedblleft{}We took the train and got off at the station.\textquotedblright{}
  \item \textbf{\ja{それから、すこし} \ja{あるきました。}} (\emph{sorekara, sukoshi arukimashita.}) — \textquotedblleft{}After that, we walked a little.\textquotedblright{}
  \item \textbf{\ja{ついた} \ja{とき、あめが} \ja{ふっていました。}} (\emph{tsuita toki, ame ga futte imashita.}) — \textquotedblleft{}When we arrived, it was raining.\textquotedblright{}
  \item \textbf{\ja{あめは} \ja{ひるに} \ja{やみました。}} (\emph{ame wa hiru ni yamimashita.}) — \textquotedblleft{}The rain stopped at noon.\textquotedblright{}
  \item \textbf{\ja{ちちと} \ja{かわで} \ja{およぎました。}} (\emph{chichi to kawa de oyogimashita.}) — \textquotedblleft{}I swam in the river with my father.\textquotedblright{}
  \item \textbf{\ja{みずは} \ja{きれいでした。}} (\emph{mizu wa kirei deshita.}) — \textquotedblleft{}The water was clean.\textquotedblright{}
  \item \textbf{\ja{ひるごはんを} \ja{たべてから、さかなを} \ja{みました。}} (\emph{hirugohan o tabete kara, sakana o mimashita.}) — \textquotedblleft{}After eating lunch, we watched the fish.\textquotedblright{}
  \item \textbf{\ja{ちちは} \ja{しんぶんを} \ja{よんでいました。}} (\emph{chichi wa shinbun o yonde imashita.}) — \textquotedblleft{}My father was reading the newspaper.\textquotedblright{}
  \item \textbf{\ja{そして、ゆうがた} \ja{うちに} \ja{かえりました。}} (\emph{soshite, yūgata uchi ni kaerimashita.}) — \textquotedblleft{}And in the evening we went home.\textquotedblright{}
  \item \textbf{\ja{かわの} \ja{みずは} \ja{年}\ja{じゅう} \ja{つめたいです。}} (\emph{kawa no mizu wa nenjū tsumetai desu.}) — \textquotedblleft{}The river water is cold all year round.\textquotedblright{}
  \item \textbf{\ja{月日}\ja{は} \ja{はやいです。}} (\emph{tsukihi wa hayai desu.}) — \textquotedblleft{}Time flies.\textquotedblright{}
\end{enumerate}

\emph{Read it:} the thirteen lines again, aloud, without the glosses
\end{tcolorbox}

\subsection*{Your turn}
Say these aloud:

\begin{itemize}
\raggedright
  \item who went to the river — \emph{chichi to}
  \item what the father was doing after lunch — \emph{shinbun o yonde imashita}
  \item when they went home — \emph{yūgata}
\end{itemize}

\begin{tcolorbox}[breakable,colback=teal!4,colframe=teal!35!black,title={Before you move on}]
Which lines set a scene, going on? (Five and ten: \textbf{\ja{ふっていました}}, \textbf{\ja{よんでいました}}.) Which word puts lunch first? (\textbf{\ja{たべてから}}.)
\end{tcolorbox}
\section[(review)]{Review — what was going on}

\label{lesson:JA-R253-what-was-going-on}



\begin{quote}
Say \textquotedblleft{}It is raining now\textquotedblright{}. (\emph{ima ame ga futte imasu}.) Then \textquotedblleft{}when I got up\textquotedblright{}. (\emph{okita toki}.)
\end{quote}

\subsection*{Your turn}
Say these aloud:

\begin{itemize}
\raggedright
  \item I am waiting for a friend.
  \item Yesterday it was raining, and the train stopped.
  \item When I was a child, I went to the river.
  \item It is cold here all year round.
  \item My father was reading the newspaper.
  \item I brush my teeth before I go to bed.
\end{itemize}

\begin{tcolorbox}[breakable,colback=teal!4,colframe=teal!35!black,title={Before you move on}]
\textquotedblleft{}Was reading\textquotedblright{}? (\textbf{\ja{よんでいました}}.) \textquotedblleft{}Is waiting\textquotedblright{}? (\textbf{\ja{まっています}}.) \textquotedblleft{}When\textquotedblright{}? (\textbf{\ja{とき}}, after the plain form.) \textquotedblleft{}All year round\textquotedblright{}? (\textbf{\ja{年}\ja{じゅう}}.)
\end{tcolorbox}
